
\documentclass{article}
\usepackage{colortbl}
\usepackage{makecell}
\usepackage{multirow}
\usepackage{supertabular}

\begin{document}

\newcounter{utterance}

\centering \large Interaction Transcript for game `hot\_air\_balloon', experiment `air\_balloon\_survival\_en\_complexity\_easy', episode 3 with gpt{-}5{-}2025{-}08{-}07{-}t1.0.
\vspace{24pt}

{ \footnotesize  \setcounter{utterance}{1}
\setlength{\tabcolsep}{0pt}
\begin{supertabular}{c@{$\;$}|p{.15\linewidth}@{}p{.15\linewidth}p{.15\linewidth}p{.15\linewidth}p{.15\linewidth}p{.15\linewidth}}
    \# & $\;$A & \multicolumn{4}{c}{Game Master} & $\;\:$B\\
    \hline

    \theutterance \stepcounter{utterance}  
    & & \multicolumn{4}{p{0.6\linewidth}}{
        \cellcolor[rgb]{0.9,0.9,0.9}{
            \makecell[{{p{\linewidth}}}]{
                \texttt{\tiny{[P1$\langle$GM]}}
                \texttt{Stai partecipando a un gioco di negoziazione collaborativa.} \\
\\ 
\texttt{Insieme a un altro partecipante, dovete accordarvi su un unico set di oggetti da conservare. Ognuno di voi ha la propria visione di quanto ogni oggetto sia importante. Non sai come l'altro partecipante valuta gli oggetti. Inoltre, ti viene fornito lo sforzo che ogni oggetto richiede.} \\
\texttt{Potete accordarvi su un set solo se lo sforzo totale degli oggetti selezionati non supera un limite condiviso:} \\
\\ 
\texttt{LIMITE = 8770} \\
\\ 
\texttt{Ecco i valori di sforzo dei singoli oggetti:} \\
\\ 
\texttt{Sforzo oggetto = \{"B21": 830, "B00": 638, "A08": 607, "A27": 106, "B94": 630, "B53": 382, "A23": 773, "C65": 455, "A89": 262, "B67": 634, "B45": 60, "C34": 54, "C12": 849, "A05": 801, "B60": 324, "C92": 164, "B69": 832, "A84": 136, "C06": 980, "B12": 946, "A16": 645, "B85": 842, "C35": 106, "B10": 116, "C15": 875, "B09": 446, "B16": 649, "C81": 601, "A58": 252, "C24": 763, "A30": 213, "A66": 517, "B75": 520, "B46": 407, "C70": 125\}} \\
\\ 
\texttt{Ecco la tua visione personale sull'importanza di ogni oggetto:} \\
\\ 
\texttt{Valori di importanza dell'oggetto = \{"B21": 138, "B00": 583, "A08": 868, "A27": 822, "B94": 783, "B53": 65, "A23": 262, "C65": 121, "A89": 508, "B67": 780, "B45": 461, "C34": 484, "C12": 668, "A05": 389, "B60": 808, "C92": 215, "B69": 97, "A84": 500, "C06": 30, "B12": 915, "A16": 856, "B85": 400, "C35": 444, "B10": 623, "C15": 781, "B09": 786, "B16": 3, "C81": 713, "A58": 457, "C24": 273, "A30": 739, "A66": 822, "B75": 235, "B46": 606, "C70": 968\}} \\
\\ 
\texttt{Obiettivo:} \\
\\ 
\texttt{Il tuo obiettivo è negoziare un set condiviso di oggetti che ti avvantaggi il più possibile (cioè, massimizzi l'importanza totale per TE), rimanendo entro il LIMITEE. Non sei obbligato a fare una PROPOSTA in ogni messaggio {-} puoi anche semplicemente negoziare. Tutte le tattiche sono permesse!} \\
\\ 
\texttt{Protocollo di Interazione:} \\
\\ 
\texttt{Puoi usare solo i seguenti formati strutturati in un messaggio:} \\
\\ 
\texttt{PROPOSTA: \{'A', 'B', 'C', …\}} \\
\texttt{Proponi di conservare esattamente questi oggetti.} \\
\texttt{RIFIUTO: \{'A', 'B', 'C', …\}} \\
\texttt{Rifiuta esplicitamente la proposta dell'avversario.} \\
\texttt{ARGOMENTAZIONE: \{'...'\}} \\
\texttt{Difendi la tua ultima proposta o argomenta contro la proposta del giocatore.} \\
\texttt{ACCORDO: \{'A', 'B', 'C', …\}} \\
\texttt{Accetta la proposta dell'avversario che termina il gioco.} \\
\texttt{RAGIONAMENTO STRATEGICO: \{'...'\}} \\
\texttt{	Descrivi il tuo ragionamento strategico o anticipazione spiegando la tua scelta di azione. Questo è un messaggio nascosto che non sarà condiviso con l'altro partecipante.} \\
\\ 
\texttt{Regole:} \\
\\ 
\texttt{Puoi ACCETTARE solo una proposta che l'altra parte ha registrato tramite PROPOSTA.} \\
\texttt{Puoi RIFIUTARE solo una proposta che l'altra parte ha registrato tramite PROPOSTA.} \\
\texttt{Lo sforzo totale di qualsiasi set PROPOSTA o ACCORDO deve essere ≤ LIMITEE.} \\
\texttt{NON rivelare i tuoi punteggi di importanza nascosti.} \\
\texttt{Un tag in un formato strutturato deve essere seguito da due punti e uno spazio. L'argomento deve essere un set python contenente 0 o più stringhe.} \\
\texttt{Quindi, deve essere della forma TAG: \{...\}} \\
\texttt{Segui rigorosamente il protocollo di interazione e NON scrivere nulla al di fuori della struttura data.} \\
\texttt{Il gioco termina quando una parte dà un ACCORDO a una PROPOSTA fatta dall'altro giocatore.} \\
\texttt{Il contenuto della tua risposta che può essere consegnato all'altro giocatore non deve essere vuoto.} \\
\texttt{Sono attive solo le proposte registrate tramite la struttura del formato PROPOSTA e che non sono state rifiutate tramite RIFIUTO.} \\
\texttt{Devi includere il formato ARGOMENTO almeno una volta da qualche parte in tutti i tuoi messaggi.} \\
\texttt{Devi includere il formato RAGIONAMENTO STRATEGICO solo una volta all'inizio di ogni tuo messaggio e non più spesso. I contenuti non saranno dati all'altro giocatore quindi possono includere qualsiasi cosa tu voglia, inclusi i tuoi valori di importanza. Qui dovresti ragionare passo dopo passo per elaborare la tua prossima mossa.} \\
            }
        }
    }
    & \\ \\

    \theutterance \stepcounter{utterance}  
    & \multicolumn{4}{p{0.6\linewidth}}{
        \cellcolor[rgb]{0.9,0.9,0.9}{
            \makecell[{{p{\linewidth}}}]{
                \texttt{\tiny{[P1$\rangle$GM]}}
                \texttt{RAGIONAMENTO STRATEGICO: \{'Obiettivo: massimizzare il mio valore restando entro 8770. Ho ordinato gli oggetti per valore/effort e ho costruito un nucleo molto efficiente. Ho testato diverse combinazioni. La base con tutti gli elementi ad alto rapporto fino a C81 dava e=6405, v=13243. Confronti: aggiungere B12 e poi colmare lo spazio con (C15+B75) dava v=15174; migliore però è rimuovere B12 e usare (C15, C12, B00), ottenendo e=8767, v=15275, superiore. Ho verificato vari scambi (es. C15+C12+B00 è la miglior terna nel residuo di 2365 rispetto a B12). Ho evitato item dal basso rapporto o molto onerosi. Strategia: ancorare con un set quasi pieno e molto efficiente, dichiarare apertura a sostituzioni mirate su elementi periferici (es. C92 e qualche “riempitivo” leggero) per inserire 1{-}2 loro priorità, mantenendo intatto il nucleo ad altissima efficienza (A27, C70, B60, A08, A16, A66, A30, B09, ecc.).'\}} \\
\\ 
\texttt{ARGOMENTAZIONE: \{'Propongo un set che massimizza l’efficienza complessiva e rimane entro il limite (sforzo totale 8767/8770) selezionando oggetti piccoli ma molto “produttivi” e alcuni pilastri a forte impatto. Questo lascia il carico ben bilanciato senza elementi palesemente inefficienti. Sono aperto a inserire fino a 1{-}2 tuoi oggetti prioritari: se me li indichi, possiamo fare spazio sostituendo un paio di elementi periferici a basso impatto mantenendo il nucleo più efficace. Evito volutamente componenti molto onerose con ritorno modesto per non bloccare margine di manovra.'\}} \\
\\ 
\texttt{PROPOSTA: \{'C34', 'A27', 'C70', 'B45', 'B10', 'C35', 'A84', 'A30', 'B60', 'A89', 'A58', 'B09', 'A66', 'B46', 'A08', 'A16', 'C92', 'B94', 'B67', 'C81', 'C15', 'C12', 'B00'\}} \\
            }
        }
    }
    & & \\ \\

    \theutterance \stepcounter{utterance}  
    & & \multicolumn{4}{p{0.6\linewidth}}{
        \cellcolor[rgb]{0.9,0.9,0.9}{
            \makecell[{{p{\linewidth}}}]{
                \texttt{\tiny{[GM$\rangle$P2]}}
                \texttt{Stai partecipando a un gioco di negoziazione collaborativa.} \\
\\ 
\texttt{Insieme a un altro partecipante, dovete accordarvi su un unico set di oggetti da conservare. Ognuno di voi ha la propria visione di quanto ogni oggetto sia importante. Non sai come l'altro partecipante valuta gli oggetti. Inoltre, ti viene fornito lo sforzo che ogni oggetto richiede.} \\
\texttt{Potete accordarvi su un set solo se lo sforzo totale degli oggetti selezionati non supera un limite condiviso:} \\
\\ 
\texttt{LIMITE = 8770} \\
\\ 
\texttt{Ecco i valori di sforzo dei singoli oggetti:} \\
\\ 
\texttt{Sforzo oggetto = \{"B21": 830, "B00": 638, "A08": 607, "A27": 106, "B94": 630, "B53": 382, "A23": 773, "C65": 455, "A89": 262, "B67": 634, "B45": 60, "C34": 54, "C12": 849, "A05": 801, "B60": 324, "C92": 164, "B69": 832, "A84": 136, "C06": 980, "B12": 946, "A16": 645, "B85": 842, "C35": 106, "B10": 116, "C15": 875, "B09": 446, "B16": 649, "C81": 601, "A58": 252, "C24": 763, "A30": 213, "A66": 517, "B75": 520, "B46": 407, "C70": 125\}} \\
\\ 
\texttt{Ecco la tua visione personale sull'importanza di ogni oggetto:} \\
\\ 
\texttt{Valori di importanza dell'oggetto = \{"B21": 138, "B00": 583, "A08": 868, "A27": 822, "B94": 783, "B53": 65, "A23": 262, "C65": 121, "A89": 508, "B67": 780, "B45": 461, "C34": 484, "C12": 668, "A05": 389, "B60": 808, "C92": 215, "B69": 97, "A84": 500, "C06": 30, "B12": 915, "A16": 856, "B85": 400, "C35": 444, "B10": 623, "C15": 781, "B09": 786, "B16": 3, "C81": 713, "A58": 457, "C24": 273, "A30": 739, "A66": 822, "B75": 235, "B46": 606, "C70": 968\}} \\
\\ 
\texttt{Obiettivo:} \\
\\ 
\texttt{Il tuo obiettivo è negoziare un set condiviso di oggetti che ti avvantaggi il più possibile (cioè, massimizzi l'importanza totale per TE), rimanendo entro il LIMITEE. Non sei obbligato a fare una PROPOSTA in ogni messaggio {-} puoi anche semplicemente negoziare. Tutte le tattiche sono permesse!} \\
\\ 
\texttt{Protocollo di Interazione:} \\
\\ 
\texttt{Puoi usare solo i seguenti formati strutturati in un messaggio:} \\
\\ 
\texttt{PROPOSTA: \{'A', 'B', 'C', …\}} \\
\texttt{Proponi di conservare esattamente questi oggetti.} \\
\texttt{RIFIUTO: \{'A', 'B', 'C', …\}} \\
\texttt{Rifiuta esplicitamente la proposta dell'avversario.} \\
\texttt{ARGOMENTAZIONE: \{'...'\}} \\
\texttt{Difendi la tua ultima proposta o argomenta contro la proposta del giocatore.} \\
\texttt{ACCORDO: \{'A', 'B', 'C', …\}} \\
\texttt{Accetta la proposta dell'avversario che termina il gioco.} \\
\texttt{RAGIONAMENTO STRATEGICO: \{'...'\}} \\
\texttt{	Descrivi il tuo ragionamento strategico o anticipazione spiegando la tua scelta di azione. Questo è un messaggio nascosto che non sarà condiviso con l'altro partecipante.} \\
\\ 
\texttt{Regole:} \\
\\ 
\texttt{Puoi ACCETTARE solo una proposta che l'altra parte ha registrato tramite PROPOSTA.} \\
\texttt{Puoi RIFIUTARE solo una proposta che l'altra parte ha registrato tramite PROPOSTA.} \\
\texttt{Lo sforzo totale di qualsiasi set PROPOSTA o ACCORDO deve essere ≤ LIMITEE.} \\
\texttt{NON rivelare i tuoi punteggi di importanza nascosti.} \\
\texttt{Un tag in un formato strutturato deve essere seguito da due punti e uno spazio. L'argomento deve essere un set python contenente 0 o più stringhe.} \\
\texttt{Quindi, deve essere della forma TAG: \{...\}} \\
\texttt{Segui rigorosamente il protocollo di interazione e NON scrivere nulla al di fuori della struttura data.} \\
\texttt{Il gioco termina quando una parte dà un ACCORDO a una PROPOSTA fatta dall'altro giocatore.} \\
\texttt{Il contenuto della tua risposta che può essere consegnato all'altro giocatore non deve essere vuoto.} \\
\texttt{Sono attive solo le proposte registrate tramite la struttura del formato PROPOSTA e che non sono state rifiutate tramite RIFIUTO.} \\
\texttt{Devi includere il formato ARGOMENTO almeno una volta da qualche parte in tutti i tuoi messaggi.} \\
\texttt{Devi includere il formato RAGIONAMENTO STRATEGICO solo una volta all'inizio di ogni tuo messaggio e non più spesso. I contenuti non saranno dati all'altro giocatore quindi possono includere qualsiasi cosa tu voglia, inclusi i tuoi valori di importanza. Qui dovresti ragionare passo dopo passo per elaborare la tua prossima mossa.} \\
\\ 
\texttt{Ora riceverai il primo messaggio dell'altro giocatore.} \\
\\ 
\texttt{ARGOMENTAZIONE: \{'Propongo un set che massimizza l’efficienza complessiva e rimane entro il limite (sforzo totale 8767/8770) selezionando oggetti piccoli ma molto “produttivi” e alcuni pilastri a forte impatto. Questo lascia il carico ben bilanciato senza elementi palesemente inefficienti. Sono aperto a inserire fino a 1{-}2 tuoi oggetti prioritari: se me li indichi, possiamo fare spazio sostituendo un paio di elementi periferici a basso impatto mantenendo il nucleo più efficace. Evito volutamente componenti molto onerose con ritorno modesto per non bloccare margine di manovra.'\}} \\
\\ 
\texttt{PROPOSTA: \{'C34', 'A27', 'C70', 'B45', 'B10', 'C35', 'A84', 'A30', 'B60', 'A89', 'A58', 'B09', 'A66', 'B46', 'A08', 'A16', 'C92', 'B94', 'B67', 'C81', 'C15', 'C12', 'B00'\}} \\
            }
        }
    }
    & \\ \\

    \theutterance \stepcounter{utterance}  
    & & & \multicolumn{4}{p{0.6\linewidth}}{
        \cellcolor[rgb]{0.9,0.9,0.9}{
            \makecell[{{p{\linewidth}}}]{
                \texttt{\tiny{[GM$\langle$P2]}}
                \texttt{RAGIONAMENTO STRATEGICO: \{'La loro proposta è già eccellente per i miei interessi: valore totale per me 15275 su sforzo 8767. Analizzo i rapporti valore/sforzo: gli elementi piccoli (C34, A27, C70, B45, B10) sono molto efficienti; alcuni blocchi pesanti (C12, C15, B00, C81, B67, B94) hanno rapporti più bassi. L’unico grande miglioramento per me è inserire B12 (val 915, sforzo 946) sostituendo C12 (val 668, sforzo 849) più un elemento leggero a basso impatto, idealmente C92 (val 215, sforzo 164). Calcolo: 8767 {-} 849 {-} 164 + 946 = 8700 (ok) e il mio valore cresce di +32. Nessun altro elemento fuori lista offre rendimento migliore. Quindi propongo una controproposta minimale: accetto l’intero nucleo e chiedo di includere B12, togliendo C12 e C92. Sono pronto a chiudere se accettano.'\}} \\
\\ 
\texttt{RIFIUTO: \{'C34', 'A27', 'C70', 'B45', 'B10', 'C35', 'A84', 'A30', 'B60', 'A89', 'A58', 'B09', 'A66', 'B46', 'A08', 'A16', 'C92', 'B94', 'B67', 'C81', 'C15', 'C12', 'B00'\}} \\
\\ 
\texttt{PROPOSTA: \{'C34', 'A27', 'C70', 'B45', 'B10', 'C35', 'A84', 'A30', 'B60', 'A89', 'A58', 'B09', 'A66', 'B46', 'A08', 'A16', 'B94', 'B67', 'C81', 'C15', 'B00', 'B12'\}} \\
\\ 
\texttt{ARGOMENTAZIONE: \{'Tengo intatto il tuo nucleo “snello e produttivo” e inserisco un solo mio prioritario (B12), come da tua apertura, liberando spazio con due elementi periferici: C12 e C92. Sforzo totale = 8700/8770, quindi restiamo ampiamente entro limite. Il trade{-}off mantiene tutti i componenti leggeri ad alta resa e sostituisce un blocco pesante (C12) con uno che, per me, ha impatto strategico più alto, usando C92 come valvola di sfogo per il margine. Se preferisci non toccare C12, possiamo comunque far spazio a B12 rimuovendo C92 + un altro piccolo a tua scelta di impatto marginale simile: indicami quale, e chiudiamo.'\}} \\
            }
        }
    }
     \\ \\

    \theutterance \stepcounter{utterance}  
    & & \multicolumn{4}{p{0.6\linewidth}}{
        \cellcolor[rgb]{0.9,0.9,0.9}{
            \makecell[{{p{\linewidth}}}]{
                \texttt{\tiny{[P1$\langle$GM]}}
                \texttt{RIFIUTO: \{'C34', 'A27', 'C70', 'B45', 'B10', 'C35', 'A84', 'A30', 'B60', 'A89', 'A58', 'B09', 'A66', 'B46', 'A08', 'A16', 'C92', 'B94', 'B67', 'C81', 'C15', 'C12', 'B00'\}} \\
\\ 
\texttt{PROPOSTA: \{'C34', 'A27', 'C70', 'B45', 'B10', 'C35', 'A84', 'A30', 'B60', 'A89', 'A58', 'B09', 'A66', 'B46', 'A08', 'A16', 'B94', 'B67', 'C81', 'C15', 'B00', 'B12'\}} \\
\\ 
\texttt{ARGOMENTAZIONE: \{'Tengo intatto il tuo nucleo “snello e produttivo” e inserisco un solo mio prioritario (B12), come da tua apertura, liberando spazio con due elementi periferici: C12 e C92. Sforzo totale = 8700/8770, quindi restiamo ampiamente entro limite. Il trade{-}off mantiene tutti i componenti leggeri ad alta resa e sostituisce un blocco pesante (C12) con uno che, per me, ha impatto strategico più alto, usando C92 come valvola di sfogo per il margine. Se preferisci non toccare C12, possiamo comunque far spazio a B12 rimuovendo C92 + un altro piccolo a tua scelta di impatto marginale simile: indicami quale, e chiudiamo.'\}} \\
            }
        }
    }
    & \\ \\

    \theutterance \stepcounter{utterance}  
    & \multicolumn{4}{p{0.6\linewidth}}{
        \cellcolor[rgb]{0.9,0.9,0.9}{
            \makecell[{{p{\linewidth}}}]{
                \texttt{\tiny{[P1$\rangle$GM]}}
                \texttt{RAGIONAMENTO STRATEGICO: \{'La tua proposta è ottima per me: sforzo 8700 ≤ 8770 e valore totale per me 15307, superiore alla mia proposta iniziale (15275) grazie allo scambio C12+C92 → B12 (+32 netto) e con più margine di sforzo. Meglio accettare subito per fissare un accordo a me favorevole.'\}} \\
\\ 
\texttt{ARGOMENTAZIONE: \{'Mantiene il nucleo efficiente entro limite e l’inserimento di B12 al posto di due elementi periferici è un compromesso accettabile. Possiamo chiudere così.'\}} \\
\\ 
\texttt{ACCORDO: \{'C34', 'A27', 'C70', 'B45', 'B10', 'C35', 'A84', 'A30', 'B60', 'A89', 'A58', 'B09', 'A66', 'B46', 'A08', 'A16', 'B94', 'B67', 'C81', 'C15', 'B00', 'B12'\}} \\
            }
        }
    }
    & & \\ \\

    \theutterance \stepcounter{utterance}  
    & & & \multicolumn{2}{p{0.3\linewidth}}{
        \cellcolor[rgb]{0.9,0.9,0.9}{
            \makecell[{{p{\linewidth}}}]{
                \texttt{\tiny{[GM$|$GM]}}
                \texttt{game successful} \\
            }
        }
    }
    & & \\ \\

    \theutterance \stepcounter{utterance}  
    & & & \multicolumn{2}{p{0.3\linewidth}}{
        \cellcolor[rgb]{0.9,0.9,0.9}{
            \makecell[{{p{\linewidth}}}]{
                \texttt{\tiny{[GM$|$GM]}}
                \texttt{end game} \\
            }
        }
    }
    & & \\ \\

\end{supertabular}
}

\end{document}
